\clearpage
\section{Conclusions}
\label{sec:conclusions}

This thesis examined shared--private representation learning for unsupervised multi-omics cancer subtyping, with emphasis on whether latent codes named shared ($S$) and private ($P$) actually realise that decomposition. Using TCGA BRCA multi-omics data and multi-seed evaluation over seeds $\{21,\ldots,25\}$, the work audited MOCSS and OMIDIENT, proposed Methods~1--3 that enforce clearer role separation, and ranked models by a probe-centred protocol rather than by private-subtype NMI alone.

The first conclusion is negative but central: high unsupervised subtype NMI does not imply faithful shared--private roles. In MOCSS, private concatenation NMI approximately matches full concatenation ($\approx0.47$ vs.\ $\approx0.46$) and exceeds mean shared NMI ($\approx0.42$), while $\Delta\mathrm{X}\,R^{2}$ remains near~$0.04$ and UMAP shows subtype structure in private embeddings. In OMIDIENT, near-zero private NMI and large $\Delta\mathrm{X}\,R^{2}$ ($\approx1.11$) coexist with collapsed own-view private reconstructiveness ($X\leftarrow P\approx-0.34$), so the apparent separation is not evidence of a verified shared theory. Single-seed optimism further shrinks under multi-seed aggregation.

The second conclusion is methodological. Models should be judged jointly by cross-view $\Delta\mathrm{X}\,R^{2}$, qualitative UMAP role separation, and integration NMI. Minimising $\mathrm{concat}(P)$ subtype NMI is not a valid primary objective: Method~1+$b_m$ and Method~1+VICReg achieve lower private NMI (approximately~$0.24$--$0.25$) without securing stable positive $\Delta\mathrm{X}\,R^{2}$ or competitive concatenation performance, whereas Method~3 tolerates high private NMI ($\approx0.44$) while delivering the clearest AE-private cross-view gap.

The third conclusion is constructive. Among the proposed approaches, Method~3 provides the strongest empirical support for improved shared--private distinction on BRCA ($\Delta\mathrm{X}\,R^{2}\approx0.25$ with concatenation NMI$\approx0.46$), subject to the spectral ceiling of the offline fused Laplacian target ($\mathrm{NMI}(U)\approx0.387$). Method~2 is the preferred simpler AE-private alternative to Method~1+VICReg, combining joint reconstruction, VICReg independence, and competitive multi-seed integration (concatenation NMI$\approx0.43$, $\Delta\mathrm{X}\,R^{2}\approx+0.05$) with structured but non-identical shared and private geometries.

KIRC experiments are reserved for a later revision. Future work should improve fused-graph targets for Method~3, harden predictability estimators beyond Method~1's two-step adversaries, test whether non-collapsing private priors can be combined with Method~2-style joint reconstruction, and extend the probe suite to additional tumour cohorts and omics combinations. More broadly, deep multi-omics models that advertise shared--private structure should treat that structure as an empirical claim to be tested, not as a consequence of architecture diagrams or clustering leaderboards.
